\documentclass[12pt]{article}
\usepackage{amsmath, amssymb}
\usepackage[margin=1in]{geometry}
\setlength{\parskip}{6pt}
\title{QUBO Formulation: RNA Secondary-Structure Stem Selection}
\date{}
\begin{document}
\maketitle

\subsection*{RNA Secondary-Structure Stem Selection}

\paragraph{Problem Statement.}
Given a set of $N$ candidate stems indexed by $i \in \{0, 1, \dots, N-1\}$, each associated with a positive stability score $w_i > 0$, and a set of pairwise conflict pairs $C \subseteq \{0, 1, \dots, N-1\} \times \{0, 1, \dots, N-1\}$, the goal is to select a subset of stems that maximizes the total stability score while ensuring that no two conflicting stems are selected simultaneously.

\paragraph{Decision Variables.}
For each candidate stem $i \in \{0, 1, \dots, N-1\}$, introduce a binary decision variable $x_i \in \{0, 1\}$. The variable $x_i$ takes the value $1$ if stem $i$ is selected, and $0$ otherwise.

\paragraph{QUBO Derivation.}
\textbf{Step 1.} The natural objective is to maximize the total stability score:
\begin{align}
    \max_{x} \sum_{i=0}^{N-1} w_i x_i.
\end{align}
Since QUBO solvers minimize energy, we negate the objective to obtain a minimization problem:
\begin{align}
    \min_{x} \left( -\sum_{i=0}^{N-1} w_i x_i \right).
\end{align}

\textbf{Step 2.} The compatibility constraints require that for every conflicting pair $(i,j) \in C$, the stems cannot both be selected:
\begin{align}
    x_i + x_j \le 1, \quad \forall (i,j) \in C.
\end{align}

\textbf{Step 3.} Each inequality constraint is converted to a quadratic penalty term. For a pair $(i,j) \in C$, the squared penalty is:
\begin{align}
    \lambda (x_i + x_j - 1)^2.
\end{align}
Expanding the square yields:
\begin{align}
    \lambda (x_i^2 + x_j^2 + 2x_i x_j - 2x_i - 2x_j + 1).
\end{align}
Applying the idempotent property $x_k^2 = x_k$ for binary variables gives:
\begin{align}
    \lambda (x_i + x_j + 2x_i x_j - 2x_i - 2x_j + 1) = \lambda (1 - x_i - x_j + 2x_i x_j).
\end{align}
Constant terms do not affect the minimizer, and linear terms can be absorbed into the objective. The essential quadratic contribution that enforces exclusivity on $\{0,1\}^N$ is $\lambda x_i x_j$. Summing over all conflicts, the total penalty contribution to the Hamiltonian is:
\begin{align}
    \lambda \sum_{(i,j) \in C} x_i x_j.
\end{align}

\textbf{Step 4.} Combining the negated objective and the penalty terms yields the unconstrained QUBO Hamiltonian $H(x)$:
\begin{align}
    H(x) = -\sum_{i=0}^{N-1} w_i x_i + \lambda \sum_{(i,j) \in C} x_i x_j.
\end{align}

\textbf{Step 5.} Reading off the coefficients from $H(x) = \sum_{i,j} Q_{ij} x_i x_j + c$, the Q matrix entries and offset are given in closed form as:
\begin{align}
    Q_{i,i} &= -w_i, & i \in \{0, \dots, N-1\}, \\
    Q_{i,j} &= \lambda, & (i,j) \in C, \\
    Q_{j,i} &= \lambda, & (i,j) \in C, \\
    Q_{i,j} &= 0, & \text{otherwise}, \\
    c &= 0.
\end{align}

\paragraph{Penalty Weight Justification.}
The penalty weight $\lambda$ must satisfy $\lambda > \max_{i \in \{0, \dots, N-1\}} w_i$. This condition guarantees that the energy cost of violating any single conflict constraint (which is at least $\lambda$) strictly exceeds the maximum possible gain in the objective function from selecting any individual stem (which is at most $\max_i w_i$). Consequently, the optimizer will never prefer an infeasible assignment over a feasible one.

\paragraph{Correctness Argument.}
Minimizing $H(x)$ over $\{0,1\}^N$ recovers the optimal solution because (i) for any feasible assignment satisfying all constraints in $C$, every penalty term vanishes, reducing $H(x)$ exactly to the negated original objective; thus, minimizing $H$ is equivalent to maximizing the total stability score. (ii) For any infeasible assignment, at least one conflict constraint is violated, incurring a penalty of at least $\lambda$. By the choice of $\lambda$, this penalty strictly dominates the maximum achievable objective value. Therefore, the global minimum of $H(x)$ must correspond to a feasible assignment that maximizes the original objective.

\paragraph{Illustrative Example.}
Consider a small instance with $N=3$ candidate stems. Let the stability scores be $w_0 = 3$, $w_1 = 5$, and $w_2 = 4$. The conflict set is $C = \{(0,1), (1,2)\}$. Following the general formulation, we choose $\lambda > \max\{3, 5, 4\}$, so we set $\lambda = 6$.

Instantiating the general Hamiltonian yields:
\begin{align}
    H(x) = -3x_0 - 5x_1 - 4x_2 + 6x_0 x_1 + 6x_1 x_2.
\end{align}
The corresponding Q matrix entries are:
\begin{align}
    Q_{0,0} &= -3, & Q_{1,1} &= -5, & Q_{2,2} &= -4, \\
    Q_{0,1} &= Q_{1,0} = 6, & Q_{1,2} &= Q_{2,1} = 6, \\
    Q_{i,j} &= 0 & \text{for all other } (i,j), & c &= 0.
\end{align}
Evaluating $H(x)$ over all $2^3 = 8$ binary assignments:
\begin{align}
    \begin{array}{c|c|c}
    x & \text{Feasible?} & H(x) \\ \hline
    (0,0,0) & \text{Yes} & 0 \\
    (1,0,0) & \text{Yes} & -3 \\
    (0,0,1) & \text{Yes} & -4 \\
    (1,0,1) & \text{Yes} & -7 \\
    (0,1,0) & \text{No} & -5 \\
    (1,1,0) & \text{No} & -2 \\
    (0,1,1) & \text{No} & -3 \\
    (1,1,1) & \text{No} & 0
    \end{array}
\end{align}
The global minimum is $H(x^*) = -7$ at $x^* = (1,0,1)$. This corresponds to selecting stems $0$ and $2$, which yields a maximum total stability score of $3 + 4 = 7$. The penalty terms correctly suppress all infeasible configurations, confirming the formulation.

\end{document}